\section{Experimental Setup}
\label{sec:experimental_setup}

\subsection{Scenario and Network Model}
\label{sec:scenario_setup}

All experiments were conducted on the Hamburg vertiport network considered in this work. 
Vertiports are represented as graph nodes corresponding to averaged pad coordinates, 
and air corridors are modeled as undirected edges forming a connected graph. 
Each corridor operates at a fixed absolute cruise layer of \SI{200}{m} 
above mean sea level (MSL). This fixed-altitude, obstacle-free representation is used explicitly as a benchmark abstraction for reproducible computational evaluation. Terrain elevation and obstacle variation are therefore not modeled in this study, and the corridor layer is defined in an MSL reference to ensure comparability across experiments. Edge weights are defined by the geodesic (Haversine) distance between connected vertiports.

This abstraction enables reproducible shortest-path routing while preserving 
geometrical realism at the network level.

\subsection{System Configuration}
\label{sec:system_configuration}

All experiments were executed on a workstation equipped with an 
Intel Core i7-8650U CPU @ 1.90\,GHz (4 cores, 8 threads) and 16\,GB RAM 
under a 64-bit operating system (x64 architecture).

\begin{itemize}
    \item \textbf{Operating system:} Ubuntu 24.04 LTS (Noble Numbat).
    \item \textbf{Python environment:} Project targets Python~3.14 or newer.
    \item \textbf{C++ compiler:} \texttt{g++}~13.3.0.
    \item \textbf{CFMS version:} CFMS project version~1.0.
    \item \textbf{NDMap version:} No standalone package release is declared; internal header documentation indicates version~2.00.
\end{itemize}

The Planner (Python implementation), CFMS trajectory generator (C++), and NDMap conflict detection module were executed locally on the same machine. No distributed deployment, containerization, or artificial network latency was introduced. Consequently, reported runtimes represent baseline computational performance under single-machine execution. The software follows service-oriented separation of concerns at the code and interface level; however, the evaluated implementation does not constitute a fully distributed multi-host deployment. In particular, trajectory generation and conflict detection are invoked as local native processes, and intermediate artifacts are exchanged through local filesystem paths. Therefore, while the code structure supports modular decomposition, actual horizontal scaling across multiple hosts would require additional deployment mechanisms such as shared state management, coordinated conflict-detection data access, and inter-service orchestration.

\subsection{Planner and Generator Configuration}
\label{sec:planner_configuration}

\begin{itemize}
    \item \textbf{Global planner:} A* search on the undirected corridor graph, 
    using Haversine distance both as edge cost and as admissible heuristic 
    \(f(n) = g(n) + h(n)\).
    
    \item \textbf{Cruise altitude:} Fixed absolute cruise layer of \SI{200}{m} (MSL reference), independent of local pad elevation, terrain variation, and obstacle constraints; used as a benchmark assumption rather than an operational airspace model.
    
    \item \textbf{Trajectory generation:} CFMS-based 4D synthesis with 
    phase-dependent nominal speeds derived from vehicle performance limits 
    scaled by a safety factor.
    
    \item \textbf{Sampling:} Generated trajectories were uniformly subsampled 
    to a fixed maximum number of points to reduce payload size.
\end{itemize}

\subsection{Evaluation Methodology}
\label{sec:evaluation_methodology}

To evaluate baseline computational performance, batches of trajectory requests 
were processed sequentially. For scalability analysis, 100 independent 
origin–destination vertiport pairs were randomly sampled from the corridor graph. 
Each trajectory request was fully completed before the next was initiated, 
ensuring that measured runtimes were not influenced by parallel processing effects.

In the baseline scalability experiment, conflict detection and resolution 
mechanisms were disabled to isolate the computational cost of global routing 
and trajectory generation. Runtime measurements were recorded separately for:

\begin{itemize}
    \item A* path computation,
    \item 4D trajectory generation, and
    \item total end-to-end processing time per request.
\end{itemize}

These measurements provide insight into the relative computational 
contributions of routing and trajectory synthesis within the proposed framework.
